%%%%%     PACKS     %%%%%
\documentclass[12pt]{article}
\usepackage[margin=1in,headsep=.60in]{geometry}
\usepackage[utf8]{inputenc}
\usepackage[table, dvipsnames]{xcolor}
\usepackage{array}
\usepackage{amsmath}
\usepackage{booktabs}
\usepackage{mdframed} %For box around text
\usepackage{hyperref} %For hyperlinks in the PDF
\hypersetup{
    colorlinks=true,
    linkcolor=blue,
    filecolor=magenta,
    urlcolor=cyan,
}
\urlstyle{same}
\usepackage{amsthm}
\usepackage{amssymb}
\usepackage{amsfonts}
\usepackage{siunitx}
\usepackage{graphicx}
\usepackage{caption}
\usepackage{pgfplots}
\graphicspath{{Images/}}
\usepackage[colorinlistoftodos]{todonotes}
\usepackage{cleveref}
\usepackage[labelformat=simple]{subcaption}
\usepackage{grffile}
\usepackage{gensymb}
\usepackage{float}
\usepackage[shortlabels]{enumitem}
\usepackage{enumitem}
\setlistdepth{9}
\usepackage{biblatex} %Imports biblatex package
\addbibresource{references.bib}
\usepackage{outlines}
\usepackage{minted}
\usepackage{pdflscape}
\usepackage{everypage}
\usepackage{multirow}
\usepackage{multicol}
\usepackage{afterpage}
\usepackage[labelfont=bf, font=small]{caption}
\usepackage{upgreek}
\usepackage{tikz}
\usepackage{xcolor}
\usepackage{listings}

% style sheet for code
\definecolor{codegreen}{rgb}{0,0.6,0}
\definecolor{codegray}{rgb}{0.5,0.5,0.5}
\definecolor{codepurple}{rgb}{0.58,0,0.82}
\definecolor{backcolour}{rgb}{0.95,0.95,0.92}

\lstdefinestyle{mystyle}{
    backgroundcolor=\color{backcolour},   
    commentstyle=\color{codegreen},
    keywordstyle=\color{magenta},
    numberstyle=\tiny\color{codegray},
    stringstyle=\color{codepurple},
    basicstyle=\ttfamily\footnotesize,
    breakatwhitespace=false,         
    breaklines=true,                 
    captionpos=b,                    
    keepspaces=true,                 
    numbers=left,                    
    numbersep=5pt,                  
    showspaces=false,                
    showstringspaces=false,
    showtabs=false,                  
    tabsize=2
}

\lstset{style=mystyle}

\usepackage{lastpage}
%%%%%     COMMANDS     %%%%%
%%%% Blue box for subsection text (not figures) %%%%
\newenvironment{bluebox}
  {\begin{mdframed}[backgroundcolor=blue!5,linecolor=blue!40,roundcorner=8pt]}
  {\end{mdframed}}

%%%% Simple neutral box for figures %%%%
\newenvironment{figbox}
  {\begin{mdframed}[roundcorner=8pt,shadow=true,shadowsize=4pt,shadowcolor=black!40]}
  {\end{mdframed}}

\begin{document}

\begin{titlepage}
\newcommand{\HRule}{\rule{\linewidth}{0.5mm}}

\center

\textsc{\LARGE Aarhus University}\\[1.5cm]
\textsc{\Large Department Of Mathematics}\\[0.5cm]
\textsc{\large Introduction to Algebra 1}\\[0.5cm]
    

\HRule\\[0.4cm]
    \center 
    {\huge\bfseries  Handin 3}\\[0.4cm] % Title of your document
\HRule\\[1.5cm]

\begin{minipage}{0.4\textwidth}
        \begin{flushleft}
            \large
            \textit{Author}\\   
            Søren M. \textsc{Damsgaard}\\
                 % Your name
        \end{flushleft}
    \end{minipage}
~
    \begin{minipage}{0.4\textwidth}
        \begin{flushright}
            \large
            \textit{Student number}\\
            \textbf{202309814}\\
                 % Studienummer\
            
        \end{flushright}
    \end{minipage}

\vfill\vfill\vfill % Position the date 3/4 down the remaining page
    
    {\large\today}

\vfill\vfill
    \includegraphics[width=0.2\textwidth]{Aarhus_University_seal.png}\\[1cm] % Include a department/university logo - this will require the graphicx package
    

\vfill
\end{titlepage}
%%%%%      CHAPTERS      %%%%%
\hspace{0.02cm}
\section*{Goofy ah Math Number 3}
\#MoreConsistency
\newline
Let $a$ and $b$ be natural numbers and  $M \in \mathbb{N}$ be a common multiple of $a$ and $b$.
\subsection*{Show that the product $ab$ divides the product $Ma$ and $Mb$.}
From definition 3.8~\cite{Algebra} it holds that $M$ is an integer that is divisible by both $a$ and $b$,
hence $M$ can be defined as $a | M$ and $b | M$ and further be define per the start of Chapte 3~\cite{Algebra} as:
\begin{equation}\label{1}
 M = a \cdot q_1 \quad M = b \cdot q_2    
\end{equation}
For $a|M$, multiply both sides by $b$ and substitute $M$ with \eqref{1}
\begin{align}\label{M}
    &a|M \\
    &= ab|Mb \\
    &= ab| (a\cdot q_1)b \\
    &= ab | (ab)\cdot q_1
\end{align}
which should hold since $q_1$ is a valid integer that is not zero, hence it holds that $ab|Mb$, the same can be shown for
$ab|Ma$ where  it start woth muliplying by $a$ and then substitute with $M = b \cdot q_2$ instead. 


\subsection*{Show that the product $ab$ divides the product $\gcd(a,b)\cdot M$.}
Using bezouts identity theorem 3.3~\cite{Algebra} let 
\begin{equation}\label{bezouts2}
    \gcd(a,b) = d = ax + by
\end{equation}
take the product $\gcd(a,b)\cdot M$ and substitute using \eqref{bezouts2}
\[
(ax + by) \cdot M = (Ma)x + (Mb)y
\]
Then from earlier, substitute $Ma$ and $Mb$ with result from \eqref{M} to get
\begin{equation}\label{Mgcd}
    (Ma)x + (Mb)y = (ab\cdot q_1)x + (ab\cdot q_2)y = ab(q_1x + q_2y) 
\end{equation}
and since $(q_1x + q_2y) \in \mathbb{Z}$ per the definition of division \eqref{Mgcd} should atleast
be divisible by $ab$ hence $ab|\gcd(a,b)\cdot M$

\subsection*{Show that $\frac{ab}{\gcd(a,b)}$ divides $M$.}
It was just shown that $ab|\gcd(a,b) \cdot M$ which implies $\gcd(a,b) \cdot M = (ab)\cdot q_3$ and since $\gcd(a,b) > 0$ and it's a whole number, it can be used to divide both sides to get:
\[
M = \frac{ab}{\gcd(a,b)} \cdot q_3 
\]
which by the definition in the start of chapter 3~\cite{Algebra} shows that $\frac{ab}{\gcd(a,b)}|M$

\subsection*{Show that $\frac{ab}{\gcd(a,b)}$ is the least common multiple of $a$ and $b$.}
Let $M' = \frac{ab}{\gcd(a,b)}$ and reiterate that $M \in \mathbb{N}$ is a common multiple of $a$ and $b$.\\
To show that $M' = lcm(a,b)$ from definition 3.8~\cite{Algebra} then it has to be shown that $a|M'$ and $b|M'$ which is to show $M' = a \cdot q_4$ and $M' = b \cdot q_5$ for some whole numbers $q_i \in \mathbb{W}$.
\newline
Take $M' = \frac{ab}{\gcd(a,b)}$ and separate $a$ from the fraction and let $\frac{b}{\gcd(a,b)} = q_4$  this means we get the expression
\[
M' = a \cdot q_4
\]
which means $a|M'$, the same steps can be taken for $b$, so now it holds for both $a$ and $b$ that $a|M'$ and $b|M'$ so $M'$ must be a common multiplier of $a$ and $b$.
\newline
Now to show that $M'$ is the least common multiple it should be enough to show that $M' \leq M$ since if $M' \leq M$ then $M'$ is smaller than ANY other common multiple to $a$ and $b$ hence it must then be the least common multiple.
\newline
This should be trivial since it's already known that $\frac{ab}{\gcd(a,b)}|M = M'|M$ and since both $M',M > 0$ then it is implied by $M'|M$ that atleast $M' \leq M$.\\
\newline
In conclusion since $M$ is any common multiple of $a,b$ and $M' \leq M$ then
\[
M' = lcm(a,b) = \frac{ab}{\gcd(a,b)}
\] 
\printbibliography
\end{document}